\section{Python identification: the question determines the reusable capability}
\label{exec18:pyid-question}

A lending analyst may ask whether changing an offer improves average repayment,
or whether a particular customer who failed to repay would have succeeded under
another offer. These questions share variables but require different information.
Let $A$ be a binary intervention and $Y(a)$ the repayment outcome under its
assignment to $a$. An average contrast $E\{Y(1)-Y(0)\}$ uses two marginal
intervention distributions. The retrospective question
$P\{Y(1)=1\mid A=0,Y=0\}$ also links the factual outcome to that person's
unobserved outcome under another assignment. A package capable of computing the
first quantity does not thereby identify the second. Nor does a representation
of a counterfactual variable establish an algorithm for identifying its law.
The distinction changes the proposed product boundary: a request must state its
query class and available information before an engine is selected.
\citep{exec13_y0_dsl,exec18_pyid_contract}

\subsection{Four complementary source boundaries}
\label{exec18:pyid-packages}

The inspected DoWhy revision separates model specification, identification and
estimation. Its identification method chooses between an ID identifier and an
automatic identifier. It can also construct a graph from supplied inputs, so
the proposed banking route requires an explicit reviewed graph instead of
treating that convenience as elicitation evidence. In the selected revision,
the function assigns the \texttt{proceed\_when\_unidentifiable} option but does
not forward it to either identifier; validation must inspect the returned
estimand and assumptions. This is a static source finding, not a demonstrated
runtime failure. Optional reuse of the pinned identification boundary remains
the preferred initial path; the GCM counterfactual interfaces require a separate
review of their structural assumptions and factual-noise reconstruction.
\citep{exec13_dowhy}

The selected pgmpy adjustment interface offers a narrower complementary role:
generation of candidate sets for a specified exposure and outcome. Candidate
generation must be followed by independent causal validation. The reviewed
implementation checks exposure parents in a way that can accept a supplied
mediator set when the exposure has no parents; its \texttt{minimal\_variance}
branch is advertised but raises \texttt{NotImplementedError}. An adapter must
exclude that branch and validate the chosen total-effect adjustment criterion.
Neither a minimal separator nor a successful role check guarantees a globally
minimum-variance adjustment set. These restrictions do not establish anything
about uninspected pgmpy modules.\citep{exec13_pgmpy_reuse_primary}

The pinned y0 implementation gives a symbolic ADMG identification route.
Its conditional algorithm invokes a rule-2 exchange test and otherwise
normalizes an identified joint distribution. The inspected separation path has
the previously developed bidirected-collider counterexample: conditioning on
the collider can give probability $9/10$, whereas replacing the relevant
variable by intervention gives $1/2$. The proposed native conditional path
therefore remains disabled. Identifying the joint distribution and then
normalizing is a possible restricted alternative when that joint is identified
and its denominator is positive. This sufficient route is not a completeness
claim for conditional identification: failure to identify the entire joint
does not prove failure to identify a particular conditional target.
\citep{exec13_y0_id,exec13_y0_idc,exec13_y0_rule2,exec13_y0_separation,exec13_y0_oracles}

Ananke adds a differently organized source reference. The selected authentic
\texttt{ananke-causal} 0.5.0 distribution was uploaded on 5 December 2023; it is
not presented as the newest release. Its \texttt{OneLineID} object fixes
intervention vertices, finds relevant outcome ancestors and checks whether
each district's kernel can be obtained by fixing other vertices. The \texttt{id()}
method returns a Boolean. The \texttt{functional()} method returns a string
describing the fixing expression and prints fixing-order diagnostics. That
string is not yet a typed portable expression, a fitted estimator or an
uncertainty statement. Proposed reuse would preserve variable identities,
sums, products and fixing operations in a restricted expression representation;
it would never evaluate arbitrary returned text as program code.
\citep{exec18_pyid_one_line,exec18_pyid_contract}

Ananke's experimental identification variants also show why exception names
cannot be the scientific result contract. \texttt{OneLineAID} raises
\texttt{NotIdentifiedError} when the experimental graphs fail its ancestral
precondition; \texttt{OneLineGID} uses the same exception for nonconforming
experimental graphs. A precondition failure must be distinguished from a
valid query proved nonidentified in a declared model class. The source's
supplied-model witness routines compare observational marginals and
interventional distributions of two supplied models. They do not establish
automatic witness construction for every query or arbitrary cross-world
identification.\citep{exec18_pyid_one_line}

Other Ananke components answer different questions. Its oracle intervenes in a
supplied true model and explicitly states that identification is not considered.
Conditional normalization inside that oracle cannot be cited as identification
from observed data. The missing-data routine addresses full-law identification
with explicit \texttt{R\_} and \texttt{X\_} naming and colluding-path checks.
The \texttt{CausalEffect} estimation interface targets the average contrast
$E\{Y(1)-Y(0)\}$ and selects among adjustment-, primal- and nested-fixability
strategies. A file named \texttt{counterfactual\_mean.py} is therefore evidence
about that average-effect interface, not a guarantee for a customer's joint
potential outcomes.\citep{exec18_pyid_oracle,exec18_pyid_missing,exec18_pyid_mean}

\subsection{What even a randomized study leaves undetermined}
\label{exec18:pyid-coupling}

Consider an exact illustrative population, rather than a fitted bank dataset.
Let $A$ and a latent response type $U$ be independent fair binary variables,
and impose consistency $Y=Y(A)$ and no interference. In the stable model
$Y(0)=U$ and $Y(1)=U$. In the switching model $Y(0)=U$ and $Y(1)=1-U$.
Both are compatible with an allowed $A\to Y$ causal model when faithfulness
is not imposed. The first has no individual response to assignment; in the
second, improvements and harms cancel in the mean. Direct enumeration gives
\begin{equation}
\begin{aligned}
 P(A=a,Y=y)&=\tfrac14 &&(a,y\in\{0,1\}),\\
 P\{Y(0)=1\}&=P\{Y(1)=1\}=\tfrac12,\\
 E\{Y(1)-Y(0)\}&=0.
\end{aligned}
\label{exec18:pyid-equal-laws}
\end{equation}
Thus an arbitrarily large randomized study can recover the marginal laws in
this example without identifying their within-person coupling. Among factual
untreated failures, $A=0,Y=0$ implies $U=0$. The stable model then has
$Y(1)=0$, whereas the switching model has $Y(1)=1$. Consequently,
\begin{equation}
 P\{Y(1)=1\mid A=0,Y=0\}=
 \begin{cases}0,&\text{stable model},\\1,&\text{switching model}.\end{cases}
\label{exec18:pyid-conditional}
\end{equation}
Both conditioning events have probability $1/4$, so this disagreement is not
created by conditioning on an impossible event. Exact rational enumeration
retains the complete observed and response-type tables for both models.
\citep{exec18_pyid_coupling}

The ambiguity can be described sharply. Write $p_{ij}=P\{Y(0)=i,Y(1)=j\}$
and let $t=p_{01}$ be the proportion helped by the intervention. The two
marginal constraints and normalization force
\begin{equation}
 (p_{00},p_{01},p_{10},p_{11})
   =\left(\tfrac12-t,t,t,\tfrac12-t\right),
 \qquad 0\leq t\leq\tfrac12.
\label{exec18:pyid-sharp}
\end{equation}
Nonnegativity gives the necessary interval. Conversely, for every $t$ in that
interval, take a latent variable with these four response-type probabilities
and randomize $A$ independently. The resulting structural model attains the
same observed law and intervention marginals. This construction proves
attainability of every interior point as well as both endpoints. Therefore
$P\{Y(1)=1,Y(0)=0\}$ has sharp bounds $[0,1/2]$, and randomization yields
\begin{equation}
 P\{Y(1)=1\mid A=0,Y=0\}
 =\frac{p_{01}}{P\{Y(0)=0\}}=2t\in[0,1].
\label{exec18:pyid-bounds}
\end{equation}
This is an identified set under a stated information model, not a confidence
interval. Sampling uncertainty would require a further procedure that accounts
for estimation of the marginal restrictions.\citep{exec18_pyid_coupling}

Suppose, solely to interpret the target, a successful recovery is worth
USD400 and changing the original intervention costs USD50. The retrospective
conditional net amount is $400(2t)-50$, with sharp range from minus USD50 to
USD350. The two models reverse the sign of this accounting despite identical
average causal effects. This quantity does not establish an implementable
prospective policy: the factual failure used to select the group was not
known before the original assignment. A future servicing decision would need
its own time-indexed action, available history and outcome horizon. The
example motivates a refusal or an explicit identified set, not hindsight-based
deployment.\citep{exec18_pyid_coupling}

Additional scientific restrictions can change the answer. Monotonicity
$Y(1)\geq Y(0)$ rules out the harmed response type, so $p_{10}=t=0$ under
these equal marginals. It then forces the stable coupling. Randomization alone
does not supply that restriction. Repeated observations, mechanistic evidence
or a richer structural model would have to justify their own connection across
worlds and time; simply collecting more independent assignments does not reveal
both incompatible assignments for the same occasion. The proposed service
must record the added assumption and its scientific owner before narrowing
the set.\citep{exec18_pyid_coupling,exec18_pyid_contract}

\subsection{Integration choices follow semantics and dependency constraints}
\label{exec18:pyid-integration}

The reuse proposal is deliberately component-specific. DoWhy supplies a
candidate explicit-graph identification boundary. pgmpy supplies candidate
adjustment sets with independent admissibility checks. y0 supplies a candidate
symbolic ID route, with its native conditional path excluded pending a
separation repair. Ananke supplies a district-fixing reference and a possible
isolated comparator. Directly adopting all four packages would not create one
uniform counterfactual service. Porting their source would transfer symbolic,
graph and maintenance obligations to this project without a demonstrated
deployment advantage; it is not selected initially.
\citep{exec13_dowhy,exec13_pgmpy_reuse_primary,exec13_y0_id,exec18_pyid_one_line,exec18_pyid_contract}

There is also a concrete environment conflict. The selected Ananke release
declares NumPy at least1.23 and below2, statsmodels at least0.13.2 and below0.14,
and pgmpy at least0.1.21 and below0.2. The selected DoWhy revision requires
NumPy above2 and statsmodels at least0.14; the selected pgmpy revision declares
NumPy at least2 and statsmodels at least0.14.5. These declared intervals cannot
all be satisfied in one unchanged environment. Separate pinned processes
communicating through a named-variable result contract are a design option,
not a tested environment resolution. Ananke also imports Sympy in the
identification source without explicitly declaring it in the inspected wheel
metadata; whether a resolved transitive dependency supplies it remains to be
checked. No import failure has been observed here.
\citep{exec18_pyid_metadata,exec18_pyid_one_line,exec13_dowhy_project,exec13_pgmpy_reuse_primary}

The proposed boundary validates cycles, node identity, edge type, latent status
and experimental preconditions before dispatch. Ananke's graph constructor
contains segregation and cycle checks, but their presence and the nearby
correctness TODO do not constitute either a qualification or a demonstrated
defect. Engine diagnostics are retained separately from scientific disposition.
An identified expression, an identified set, a model-specific nonidentification
result, an unsupported question, invalid input and an engine failure remain
distinct outcomes. An empty valid adjustment set must not become a failed query,
and an exception must not become a proved impossibility.
\citep{exec18_pyid_graph,exec18_pyid_one_line,exec18_pyid_contract}

Rights are similarly tied to the actual component and distribution. The
selected DoWhy source carries MIT text and y0 carries BSD three-clause text.
The Ananke archive contains two byte-identical Apache2.0 license files, while
its metadata also includes an Other/Proprietary classifier. The conflicting
declarations remain visible for distribution review rather than being silently
resolved by a package badge. Notice retention, modified-file marking where
applicable and dependency-specific obligations belong in the selected
distribution design. Process separation does not settle those questions.
\citep{exec13_dowhy_license,exec13_y0_license,exec18_pyid_license,exec18_pyid_metadata}

This comparison leaves explicit survey work. DoWhy's GCM counterfactual paths,
implementations of general counterfactual identification and Python
partial-identification optimizers need their own authenticated technical
sources, model classes and output checks. A solver would need to distinguish
a feasible witness from a certified sharp endpoint; a counterfactual engine
would need to state what structural information permits abduction and coupling.
The four inspected families support the bounded reuse choices described here.
They do not establish exhaustive Python coverage, installed adapters or
qualified banking decisions.\citep{exec18_pyid_contract}
